% Sections 13.1-13.6, from lectures/L13/appendix/a_crc15.md.

\section{Interface}
\label{c13:sec:iface}

\modulefig[0.62]{lectures/L13/appendix/images/crc15.png}

\begin{booktable}{@{}p{5.6em}p{3.4em}p{5.2em}L@{}}
  \thead{Port} & \thead{Dir.} & \thead{Type} & \thead{Meaning} \\
  \midrule
  \code{clock} & in & \code{std_logic} & 50 MHz system clock. \\
  \code{reset_s2_n} & in & \code{std_logic} & Active-low, synchronised reset. \\
  \code{clear} & in & \code{std_logic} & Returns the register to all zeros at this edge, without a
    reset. \\
  \code{enable} & in & \code{std_logic} & Integrate \code{data} into the register at this edge when
    \code{'1'}. \\
  \code{data} & in & \code{std_logic} & The current bit, the bit being sent (TX) or just sampled from
    the bus (RX). \\
  \code{crc} & out & \code{crc_t} & The register's current 15-bit CRC value. \\
  \code{valid} & out & \code{std_logic} & \code{'1'} whenever the register currently reads all
    zeros. \\
\end{booktable}

No generics. \code{crc15_tb} binds its ports \textbf{positionally}, so the order and the types above
are what has to match; the names are yours. \code{crc_t} (a 15-bit vector), its width
\code{CRC_WIDTH} and the polynomial \code{CRC_POLY} come from \code{can_def} (\chapref{c10:ch}).

% ------------------------------------------------------------------------------------------
\section{The recursion}
\label{c13:sec:recursion}

Keep a 15-bit register (\code{crc_t}), cleared to all zeros by \code{reset_s2_n}. At every clock edge
where \code{enable = '1'}, integrate one bit:

\begin{textdrawing}
feedback = crc(CRC_WIDTH-1) xor data      -- Incoming bit XOR the register's top bit.
shifted  = crc shifted one step left, new lowest bit '0'
crc      = shifted xor CRC_POLY           -- if feedback = '1'
         = shifted                        -- otherwise
\end{textdrawing}

That is the whole engine, fifteen flip-flops and a shift XOR-gated by the polynomial's taps. When
\code{enable} is \code{'0'}, hold the register unchanged. Drive \code{crc} continuously from the
register, and \code{valid} as ``the register is all zeros''.

\begin{warning}
\code{clear} does to the register what \code{reset_s2_n} does, but it is a \textbf{synchronous
command, not a reset}: it belongs inside the clocked branch, tested before \code{enable} so that it
wins when both are high at the same edge. Do not fold it into the asynchronous reset condition; that
would make \code{clear} asynchronous and level-sensitive, a different design that happens to pass the
testbench (which always holds \code{clear} for a whole clock cycle). The interface table says ``at
this edge'' and means it.
\end{warning}

\code{clear} exists because the engine is reused for every frame and holds no per-frame state of its
own: \code{can_controller} pulses it once at the start of every frame (\chapref{c16:ch}). A frame
allowed to run to the end returns the register to zero by itself, through the
generate-then-check property below, so on the happy path \code{clear} changes nothing. A frame that
is \emph{aborted} halfway does not, and that is the case it exists for.

% ------------------------------------------------------------------------------------------
\section{One engine, two jobs}
\label{c13:sec:twojobs}

\begin{itemize}
  \item \textbf{Generate:} enable the engine over SOF through the end of the data field. Then read
    \code{crc}; that is the value to send as the frame's CRC field.
  \item \textbf{Check:} enable it over SOF through the data field \emph{and then carry on}, feeding
    in the \emph{received} CRC field's own bits as well. If the frame arrived intact, that drives the
    register back to exactly zero, a property of the arithmetic, not something the engine was told to
    expect. \code{valid} reports precisely that. \code{can_controller} (\chapref{c17:ch}) checks
    \code{valid} only at that point; it needs no second, separate checking circuit.
\end{itemize}

Both jobs consume \textbf{real bits only}. Stuff bits (\chapref{c14:ch}) never enter the CRC on
either side of the link; the caller's gating of \code{enable} is what keeps them out
(\chapref{c16:ch} and \ref{c17:ch} build exactly that), and the engine itself never knows they exist.
That is why \code{crc15} needs no knowledge of frame fields or stuffing: whoever drives \code{enable}
has already taken those decisions.

Both jobs assume the register starts a frame at zero, which is exactly what \code{clear} guarantees.
Note how much rests on that: a node abandoning a frame halfway (it loses arbitration, or sees a
stuffing violation, or fails the CRC check, all things \chapref{c17:ch} does) stops feeding the
engine with a \emph{partial} value in it. Without a per-frame clear, the next frame gets that residue
as its starting value, so the node computes a CRC nobody else agrees with and computes a wrong check
on everything it receives, for every frame from then until reset. A single aborted frame would take
the node off the bus permanently.

\code{valid} reads \code{'1'} at reset as well, before any data at all; meaningless there, meaningful
only at the check point above.

% ------------------------------------------------------------------------------------------
\section{A worked example, by hand}
\label{c13:sec:worked}

Start with the simplest case. Feed a single \code{1} into a just-reset engine: \code{feedback} is
\code{0 xor 1 = 1}, \code{shifted} is still all zeros, so \code{crc} becomes \code{0 xor CRC_POLY},
the polynomial itself. That is the first thing the testbench checks, and a good sanity check on your
own code.

Now four bits, \code{1001}, MSB first. \code{top} is the register's bit 14 \emph{before} the edge,
and \code{CRC_POLY} is \code{100010110011001}:

\begin{textdrawing}
         top  data  fb   shifted          xor poly   crc
reset      -     -   -   -                       -   000000000000000
bit 1      0     1   1   000000000000000       yes   100010110011001
bit 2      1     0   1   000101100110010       yes   100111010101011
bit 3      1     0   1   001110101010110       yes   101100011001111
bit 4      1     1   0   011000110011110        no   011000110011110
\end{textdrawing}

Read one row at a time and the engine stops looking like arithmetic. Every edge does the same three
things: XOR the incoming bit with the bit falling off the top, shift left, and XOR the polynomial
back in only if that result was \code{1}.

Two rows are worth stopping at. \textbf{Bit 2} feeds in a \code{0} and triggers the XOR all the same,
because the \code{1} leaving the top is what sets \code{feedback}; the incoming bit is only half the
decision. \textbf{Bit 4} feeds in a \code{1} and does \emph{not} trigger it, because the top bit is
also \code{1} and the two cancel. That is why a CRC responds to \emph{where} a bit sits and not just
to how many ones there are, and it is the property Exercise~\ref{c13:ex:missed} uses to catch a
corruption a simple sum cannot see.

Note as well that every row's \code{shifted} column is the previous row's \code{crc} shifted one step
left, with a \code{0} coming in from below. Nothing is ever added at the low end; the register's
contents come entirely from the polynomial being XORed in over and over at different offsets.

% ------------------------------------------------------------------------------------------
\section{What the testbench pins down}
\label{c13:sec:tb}

\begin{itemize}
  \item Just after reset, \code{crc} is all zeros and \code{valid = '1'}.
  \item A single \code{'1'} bit gives \code{crc = CRC_POLY}.
  \item \code{clear} returns a loaded register to all zeros without \code{reset_s2_n} being touched.
  \item A 19-bit sequence of SOF, ID, RTR and control field gives a CRC matching a value computed
    independently in Python (not derived from this VHDL).
  \item Feeding the just-computed CRC's own bits straight back in returns the register to all zeros
    (\code{valid = '1'}), the generate-then-check property.
\end{itemize}

% ------------------------------------------------------------------------------------------
\section{Putting it into \code{can_controller}}
\label{c13:sec:wire}

The same two-part change as \chapref{c12:ch}'s, one block further in. In
\file{controller/can_controller.vhd}, add the signals under the \code{bit_timer} group:

\begin{vhdlcode}
    -- crc15.
    signal crc_clear, crc_enable, crc_data_in, crc_valid : std_logic;
    signal crc_value                                     : crc_t;
\end{vhdlcode}

and the instance under \code{bit_timer1}:

\begin{vhdlcode}
    crc1: entity work.crc15
        port map(clock, reset_s2_n, crc_clear, crc_enable, crc_data_in, crc_value, crc_valid);
\end{vhdlcode}

Note that \code{crc15} connects to nothing but the clock, the reset and its own signals: it takes its
bits from whichever shift register is active, and both of those registers are as yet unwritten. The
connection joining them is \chapref{c16:ch}'s, and it is a single expression rather than a port map.
